
% 上に凸・下に凸の2つの放物線と正方形
\TypeQuestion{}{figure_application_square_between_opposite_parabolas}
\FigProblemBlock{45mm}{fig/lesson_06_opposite_parabolas}{\TextTall{図のように、2つの放物線 $y=\dfrac{1}{2}x^2$、$y=-\dfrac{3}{2}x^2$ の間に、各辺が座標軸に平行な四角形 $\mathrm{ABCD}$ がある。2点 $\mathrm{A},\mathrm{B}$ は放物線 $y=-\dfrac{3}{2}x^2$ 上に、2点 $\mathrm{C},\mathrm{D}$ は放物線 $y=\dfrac{1}{2}x^2$ 上にあり、点 $\mathrm{B}$ の $x$ 座標は正である。次の問いに答えなさい。}

\QQ{点 $\mathrm{C}$ の $x$ 座標が $2$ であるとき、点 $\mathrm{A}$ の座標を求めなさい。}{$(-2,\ -6)$}
\QQ{点 $\mathrm{C}$ の $x$ 座標が $2$ であるとき、四角形 $\mathrm{ABCD}$ の面積を求めなさい。}{$32$}
\QQ{長方形 $\mathrm{ABCD}$ が正方形となるとき、その1辺の長さを求めなさい。}{$2$}
}{
}

% 等積変形：下に凸
\TypeQuestion{}{figure_application_equal_area_transformation}
\FigProblemBlock{45mm}{fig/lesson_06_equal_area_01}{\TextTall{図のように、放物線 $y=x^2$ 上に2点 $\mathrm{A},\mathrm{B}$ があり、それぞれの $x$ 座標は $-1,1$ である。また、$x$ 座標が $1$ より大きい放物線上の点 $\mathrm{D}$ について、$\Tri{AOB}$ と $\Tri{DOB}$ の面積は等しい。次の問いに答えなさい。}}{
\QQ{点 $\mathrm{A},\mathrm{B}$ の座標をそれぞれ求めなさい。}{$\mathrm{A}(-1,\ 1),\ \mathrm{B}(1,\ 1)$}
\QQ{直線 $\mathrm{AD}$ の式を求めなさい。}{$y=x+2$}
\QQ{$\Tri{AOB}$ の面積を求めなさい。}{$1$}
\QQ{点 $\mathrm{D}$ の座標を求めなさい。}{$(2,\ 4)$}
}

